\input{libraries/report-preamble.tex}
\title{Euclidean GCD on Unsigned 64-Bit Integers}
\begin{document}
\maketitle

\section{Task and interface}
The task is to compute the greatest common divisor of two nonnegative integers
represented by \texttt{UInt64}.  This implementation uses Euclid's remainder
algorithm.  Its intended result is $\gcd(a,b)$, with $\gcd(0,0)=0$.
All input words are admitted.  The result fits in one word.

\texttt{LeanExe.Lib.NumberTheory.gcd} takes two words and returns one word.
The implementation belongs to the number-theory library.  It has no heap
arguments and allocates no arrays.  Formal correctness proofs for this
library component are deferred.

\section{Algorithm and source}
Each iteration replaces $(x,y)$ with $(y,x\bmod y)$ while $y>0$.
The recurrence preserves the common divisors, and the nonzero second operand
decreases.  When it reaches zero, the first operand is the result.  NIST's
algorithm entry gives this recurrence~\cite{euclid}.  The zero-input convention
extends its positive-input description.

\lstinputlisting{LeanExe/Lib/NumberTheory/Gcd.lean}

The listing is read from the executable source during the document build.
The loop performs one unsigned remainder per iteration.  For $(48,18)$,
the successive pairs are $(48,18)$, $(18,12)$, $(12,6)$, and $(6,0)$.
This gives result $6$ after three iterations.

\section{Clients and execution}
The fraction client calls the component and divides both inputs by the returned
divisor when the denominator is positive.  Its result for $48/18$
is $8/3$.  A zero numerator yields $0/1$.

The polynomial-ratio client first calls checked Horner evaluation for its
numerator and denominator, then calls fraction reduction.  At $x=2$,
$(1+3x+2x^2)/(2+2x)=15/6=5/2$.  These clients are compiled together with
their imported library functions and executed in Wasmtime.

\begin{lstlisting}
tools/seminum gcd 48 18
6
tools/seminum fraction 48 18
8/3
tools/seminum ratio 2 1,3,2 2,2
5/2
\end{lstlisting}

The command-line parser accepts unsigned decimal words through $2^{64}-1$.
Invalid syntax and a zero fraction denominator produce status 2 and a message
on standard error.  Successful commands print their result on standard output.

\section{Tests and component discovery}
\texttt{node test/seminum.js} compares WASM results with native Lean and
mathematical references.  Cases include zero inputs, maximum words, consecutive
large Fibonacci numbers, fraction reduction, and polynomial ratios.  The native
GCD test also compares the component with Lean's \texttt{Nat.gcd}.

The compiler's existing annotation file records source declarations and WASM
function indexes.  The example tool joins those records with a small catalog
to identify this implementation in the compiled program.  This inventory
provides discovery information.  Proof status is recorded as deferred.

\begin{thebibliography}{9}
\bibitem{euclid} Paul E. Black, ``Euclid's algorithm,'' NIST Dictionary of
Algorithms and Data Structures, 2015.
\url{https://xlinux.nist.gov/dads/HTML/euclidalgo.html}.
The recurrence is the algorithm used here.  The component specifies zero inputs
and a fixed-width representation in addition to that reference.
\end{thebibliography}
\end{document}
